\section{Geometry of additive updates}\label{app:geometry}
\subsection{Distance to a two-sided orthogonal transformation}
Write $q=\min(m,n)$ and let $O(k)$ denote the real $k\times k$ orthogonal group. The orbit $\{\bm Q_L\bm W_0\bm Q_R^{\!\top}:\bm Q_L\in O(m),\bm Q_R\in O(n)\}$ contains exactly the matrices with the same singular values as $\bm W_0$. It includes reflections as well as rotations. A local smooth change of singular bases is generated by skew-symmetric matrices.

\begin{proof}[Proof of Proposition~\ref{prop:orthogonal-distance}]
\textbf{Step 1. Bound the distance from below.} Let $\bm X=\bm Q_L\bm W_0\bm Q_R^{\!\top}$ be any matrix in the orbit. Orthogonal multiplication preserves singular values, so $\bm X$ has singular values $\bm s_0$ and $\|\bm X\|_{\mathrm F}=\|\bm s_0\|_2$. Write $\langle\bm A,\bm B\rangle_{\mathrm F}=\operatorname{tr}(\bm A^{\!\top}\bm B)$. The singular-value trace inequality \citep{mirsky1975trace} bounds this inner product by the sum of the paired, decreasing singular values, giving
\begin{equation*}
\langle\bm W_*,\bm X\rangle_{\mathrm F}\leq\sum_{i=1}^q s_{*,i}s_{0,i}.
\end{equation*}
Expanding the squared distance therefore gives
\begin{equation*}
\begin{aligned}
\|\bm W_*-\bm X\|_{\mathrm F}^2
&=\|\bm s_*\|_2^2+\|\bm s_0\|_2^2-2\langle\bm W_*,\bm X\rangle_{\mathrm F}\\
&\geq\sum_{i=1}^q(s_{*,i}-s_{0,i})^2=\|\bm s_*-\bm s_0\|_2^2.
\end{aligned}
\end{equation*}
\textbf{Step 2. Attain the bound by aligning the singular bases.} Take full SVDs $\bm W_j=\bm U_j\bm\Sigma_j\bm V_j^{\!\top}$ for $j\in\{0,*\}$. Here $\bm U_j$ and $\bm V_j$ are square orthogonal completions of the thin bases, and $\bm\Sigma_j\in\mathbb R^{m\times n}$ is rectangular diagonal. Choose $\bm Q_L=\bm U_*\bm U_0^{\!\top}$ and $\bm Q_R=\bm V_*\bm V_0^{\!\top}$. These matrices are orthogonal and give
\begin{equation*}
\bm Q_L\bm W_0\bm Q_R^{\!\top}=\bm U_*\bm\Sigma_0\bm V_*^{\!\top}=\bm W_R.
\end{equation*}
Thus $\bm W_R$ keeps the base singular values while using the adapted singular bases. Orthogonal invariance of the Frobenius norm yields
\begin{equation*}
\|\bm W_*-\bm W_R\|_{\mathrm F}
=\|\bm\Sigma_*-\bm\Sigma_0\|_{\mathrm F}
=\|\bm s_*-\bm s_0\|_2.
\end{equation*}
This equals the lower bound from Step 1, proving the minimum. No positivity or gap assumption was used, so the result also covers rank-deficient matrices and repeated singular values.

\textbf{Step 3. Relate the distance to $d_S/d_W$.} For nonzero base and update norms, the common normalization $\|\bm s_0\|_2=\|\bm W_0\|_{\mathrm F}$ cancels, giving
\begin{equation*}
\frac{d_S}{d_W}
=\frac{\|\bm s_*-\bm s_0\|_2}{\|\bm W_*-\bm W_0\|_{\mathrm F}}
=\frac{\|\bm W_*-\bm W_R\|_{\mathrm F}}{\|\bm W_*-\bm W_0\|_{\mathrm F}}.
\end{equation*}
For multiple matrices, summing the squared residuals before taking square roots gives the same interpretation for the pooled ratio.
\end{proof}

\subsection{First-order spectral and basis motion}
\begin{proof}[Proof of Proposition~\ref{prop:first-order-motion}]
\textbf{Step 1. Express the derivative in singular coordinates.} The positive, distinct singular values allow a smooth local choice of singular vectors \citep{stewart1990}. Complete them to smooth square orthogonal bases and write $\bm W_0+\epsilon\bm E=\bm U(\epsilon)\bm\Sigma(\epsilon)\bm V(\epsilon)^{\!\top}$. A dot denotes differentiation at $\epsilon=0$. Define
\begin{equation*}
\bm C=\bm U_0^{\!\top}\bm E\bm V_0,\qquad
\bm\Omega_U=\bm U_0^{\!\top}\dot{\bm U},\qquad
\bm\Omega_V=\bm V_0^{\!\top}\dot{\bm V}.
\end{equation*}
Differentiating $\bm U^{\!\top}\bm U=\bm I$ gives $\bm\Omega_U^{\!\top}=-\bm\Omega_U$. The same argument applies to $\bm\Omega_V$. Now differentiate the SVD and multiply on the left by $\bm U_0^{\!\top}$ and on the right by $\bm V_0$ to obtain
\begin{equation*}
\bm C=\bm\Omega_U\bm\Sigma_0+\dot{\bm\Sigma}-\bm\Sigma_0\bm\Omega_V,
\end{equation*}
where $\dot{\bm\Sigma}$ is rectangular diagonal with entries $\dot s_i$. Since skew-symmetric matrices have zero diagonal, the two basis terms contribute nothing to $C_{ii}$. Hence
\begin{equation*}
\dot s_i=C_{ii}=\bm u_{0,i}^{\!\top}\bm E\bm v_{0,i},\qquad 1\leq i\leq q.
\end{equation*}

\textbf{Step 2. Show that the other coordinates describe basis motion.} Holding the singular values fixed leaves $\bm\Omega_U\bm\Sigma_0-\bm\Sigma_0\bm\Omega_V$, the tangent directions to the orbit in singular coordinates. These directions have zero diagonal. Conversely, for $i\neq j\leq q$, the equations for any pair of off-diagonal entries are
\begin{equation*}
C_{ij}=(\Omega_U)_{ij}s_{0,j}-s_{0,i}(\Omega_V)_{ij},\qquad
C_{ji}=-(\Omega_U)_{ij}s_{0,i}+s_{0,j}(\Omega_V)_{ij}.
\end{equation*}
Solving this two-variable system gives
\begin{equation*}
(\Omega_U)_{ij}=\frac{s_{0,j}C_{ij}+s_{0,i}C_{ji}}{s_{0,j}^2-s_{0,i}^2},\qquad
(\Omega_V)_{ij}=\frac{s_{0,i}C_{ij}+s_{0,j}C_{ji}}{s_{0,j}^2-s_{0,i}^2}.
\end{equation*}
The denominators are nonzero because the singular values are distinct. If $m>n$, the extra rows satisfy $C_{aj}=(\Omega_U)_{aj}s_{0,j}$ for $a>q$. If $n>m$, the extra columns satisfy $C_{ib}=-s_{0,i}(\Omega_V)_{ib}$ for $b>q$. These equations are solvable because each $s_{0,i}>0$. Thus every coordinate outside the first $q$ diagonal entries can be produced by a change of orthogonal bases.

\textbf{Step 3. Separate the normal and tangent components.} The matrices $\bm u_{0,i}\bm v_{0,i}^{\!\top}$ are Frobenius-orthonormal because their pairwise inner products are $(\bm u_{0,i}^{\!\top}\bm u_{0,j})(\bm v_{0,i}^{\!\top}\bm v_{0,j})=\delta_{ij}$. By Step 2 they span the orthogonal complement of the tangent space. The normal projection of $\bm E$ is therefore
\begin{equation*}
\bm E_N=\sum_{i=1}^q\dot s_i\bm u_{0,i}\bm v_{0,i}^{\!\top},\qquad
\bm E_T=\bm E-\bm E_N.
\end{equation*}
The remainder $\bm E_T$ is tangent and changes only the bases to first order. Orthogonality gives
\begin{equation*}
\|\bm E\|_{\mathrm F}^2=\|\bm E_N\|_{\mathrm F}^2+\|\bm E_T\|_{\mathrm F}^2
=\|\dot{\bm s}\|_2^2+\|\bm E_T\|_{\mathrm F}^2.
\end{equation*}

\textbf{Step 4. Average the first-order spectral energy.} Let $\bm z_i=\operatorname{vec}(\bm u_{0,i}\bm v_{0,i}^{\!\top})$, so $\|\bm z_i\|_2=1$ and $\dot s_i=\bm z_i^{\!\top}\operatorname{vec}(\bm E)$. The assumed second moment yields
\begin{equation*}
\mathbb E[\dot s_i^2]
=\bm z_i^{\!\top}\mathbb E[\operatorname{vec}(\bm E)\operatorname{vec}(\bm E)^{\!\top}]\bm z_i
=\frac{\|\bm z_i\|_2^2}{mn}=\frac{1}{mn}.
\end{equation*}
Summing these $q$ equal contributions gives
\begin{equation*}
\mathbb E\|\dot{\bm s}\|_2^2=\frac{q}{mn}=\frac{1}{\max(m,n)}.
\end{equation*}
Since $\|\bm E\|_{\mathrm F}=1$, Step 3 gives $\mathbb E\|\bm E_T\|_{\mathrm F}^2=1-1/\max(m,n)$ for the complementary tangent energy.

\textbf{Step 5. Pass from the derivative to $d_S/d_W$.} The base norms cancel as in Proposition~\ref{prop:orthogonal-distance}, and $\|\epsilon\bm E\|_{\mathrm F}=|\epsilon|$. Thus, for $\epsilon\neq0$,
\begin{equation*}
\left(\frac{d_S}{d_W}\right)^2=\frac{\|\bm s(\bm W_0+\epsilon\bm E)-\bm s_0\|_2^2}{\epsilon^2}.
\end{equation*}
With $\bm W_0$ fixed, positivity and the nonzero singular-value gaps give the expansion
\begin{equation*}
\bm s(\bm W_0+\epsilon\bm E)-\bm s_0=\epsilon\dot{\bm s}+O(\epsilon^2),
\end{equation*}
whose remainder is bounded uniformly over $\|\bm E\|_{\mathrm F}=1$. Substitution into the ratio therefore gives $(d_S/d_W)^2=\|\dot{\bm s}\|_2^2+O(|\epsilon|)$ with a uniform error. Taking expectations and applying Step 4 gives
\begin{equation*}
\lim_{\epsilon\to0}\mathbb E\left[\left(\frac{d_S}{d_W}\right)^2\right]
=\mathbb E\|\dot{\bm s}\|_2^2=\frac{1}{\max(m,n)}.
\end{equation*}
Taking the square root proves the stated RMS limit.
\end{proof}

\begin{corollary}[Isotropic fixed-rank perturbations]\label{cor:isotropic-fraction}
Let $\bm W_0$ satisfy Proposition~\ref{prop:first-order-motion} and let $\bm E=\bm P\bm\Lambda\bm Q^{\!\top}/\|\bm\Lambda\|_{\mathrm F}$, where $\bm\Lambda\in\mathbb R^{r\times r}$ is fixed, nonzero diagonal and $\bm P\in\mathbb R^{m\times r},\bm Q\in\mathbb R^{n\times r}$ are independent uniform orthonormal frames, with $1\leq r\leq q$. Then \Cref{eq:relative-spectral-scale} holds.
\end{corollary}
\begin{proof}
\textbf{Step 1. Check the norm.} Since $\bm P^{\!\top}\bm P=\bm Q^{\!\top}\bm Q=\bm I_r$, multiplication by these frames preserves the Frobenius norm of $\bm\Lambda$. Hence $\|\bm E\|_{\mathrm F}=1$.

\textbf{Step 2. Compute the second moment.} Uniformity and independence of the frames make the distribution of $\bm E$ unchanged by any fixed left or right orthogonal multiplication. In particular, a row sign flip shows $\mathbb E[E_{ab}E_{cd}]=0$ when $a\neq c$. When $a=c$ but $b\neq d$, a column sign flip gives the same conclusion. Row and column permutations make all $\mathbb E[E_{ab}^2]$ equal. Their sum is $\mathbb E\|\bm E\|_{\mathrm F}^2=1$, so each equals $1/(mn)$. Consequently,
\begin{equation*}
\mathbb E[\operatorname{vec}(\bm E)\operatorname{vec}(\bm E)^{\!\top}]=\frac{\bm I_{mn}}{mn}.
\end{equation*}
\textbf{Step 3. Apply Proposition~\ref{prop:first-order-motion}.} Steps 1 and 2 verify its two assumptions on $\bm E$, so both conclusions in \Cref{eq:relative-spectral-scale} follow.
\end{proof}

\begin{corollary}[Pooling matrices]\label{cor:pooled-spectral-scale}
For a finite collection of matrices satisfying the isotropic assumptions of Proposition~\ref{prop:first-order-motion}, let $\bm W_{*,j}=\bm W_{0,j}+\epsilon a_j\bm E_j$, with fixed $a_j\geq0$ and $\sum_j a_j^2>0$. The pooled displacements satisfy
\begin{equation*}
\lim_{\epsilon\to0}\mathbb E\left[\left(\frac{d_S}{d_W}\right)^2\right]
=\frac{\sum_j a_j^2/\max(m_j,n_j)}{\sum_j a_j^2}.
\end{equation*}
\end{corollary}
\begin{proof}
\textbf{Step 1. Write the pooled ratio.} Let $\bm s_{*,j}$ and $\bm s_{0,j}$ be the singular-value vectors of the adapted and base matrices. The common base-norm normalization cancels, while $\|\bm E_j\|_{\mathrm F}=1$ fixes each update norm. Therefore
\begin{equation*}
\left(\frac{d_S}{d_W}\right)^2
=\frac{\sum_j\|\bm s_{*,j}-\bm s_{0,j}\|_2^2}{\sum_j\|\epsilon a_j\bm E_j\|_{\mathrm F}^2}
=\frac{\sum_j\|\bm s_{*,j}-\bm s_{0,j}\|_2^2}{\epsilon^2\sum_j a_j^2}.
\end{equation*}
\textbf{Step 2. Expand each numerator.} Write $\dot{\bm s}_j$ for the spectral derivative at $\bm W_{0,j}$ in direction $\bm E_j$. The uniform expansion from Proposition~\ref{prop:first-order-motion} gives, for fixed $a_j$,
\begin{equation*}
\frac{\mathbb E\|\bm s_{*,j}-\bm s_{0,j}\|_2^2}{\epsilon^2}
=a_j^2\mathbb E\|\dot{\bm s}_j\|_2^2+O(|\epsilon|)
=\frac{a_j^2}{\max(m_j,n_j)}+O(|\epsilon|).
\end{equation*}
When $a_j=0$, the corresponding term is identically zero.

\textbf{Step 3. Sum and take the limit.} There are finitely many matrices and $\sum_j a_j^2$ is fixed and positive. Substituting Step 2 into Step 1 yields
\begin{equation*}
\lim_{\epsilon\to0}\mathbb E\left[\left(\frac{d_S}{d_W}\right)^2\right]
=\frac{\sum_j a_j^2/\max(m_j,n_j)}{\sum_j a_j^2}.
\end{equation*}
Only linearity of expectation is used, so independence between matrices is unnecessary. The result is an average of the individual squared-ratio scales, weighted by their update energies $a_j^2$.
\end{proof}

Under the stated isotropic assumptions, the inverse-dimension scale explains how an additive update can have small relative spectral displacement. This motivates the empirical geometry measurements, while restoration tests the functional role of the spectral changes. 
